\documentclass[11pt,a4paper]{article}
\usepackage{planos-ifmt}
\begin{document}
\Ficha{Técnicas de integração, integral definida e aplicações}{IFMT Edital 029/2024, Anexo X retificado em 30/08/2024, Matemática, tema 3.}{Curso superior de tecnologia com Cálculo I; também adaptável a formação técnica com cálculo.}{Deduzir as regras de substituição e integração por partes, selecionar estratégias e verificar resultados pela derivada.}{Regra da cadeia, derivada do produto e integrais elementares.}{Exposição matemática com problema inicial de duas integrais distintas, soluções completas e verificação.}
\Roteiro{Problema: como integrar $\int2x\cos(x^2)\,dx$ e $\int xe^x\,dx$ sem decorar respostas?}{Relembrar regra da cadeia e produto, explicitando o sentido das substituições.}{Deduzir ambas as fórmulas a partir de derivadas; resolver integral por substituição e conferir.}{Resolver por partes e calcular $\int_0^1 xe^x\,dx=1$; discutir quando as técnicas não são diretamente adequadas.}{Concluir com um quadro comparativo e uma pergunta sobre verificação de primitivas.}
\section{Demonstrações}
\begin{proposicao}[Substituição]
Se $F'=f$ e $u$ é diferenciável, então $\int f(u(x))u'(x)\,dx=F(u(x))+C$ nos intervalos de definição.
\end{proposicao}
\begin{proof}
Pela regra da cadeia, $\frac{d}{dx}F(u(x))=F'(u(x))u'(x)=f(u(x))u'(x)$.
\end{proof}
\begin{proposicao}[Integração por partes]
Para funções $u,v$ de classe $C^1$, $\int u v'\,dx=uv-\int u'v\,dx$.
\end{proposicao}
\begin{proof}
A regra do produto fornece $(uv)'=u'v+uv'$. Integre e isole o termo $\int uv'$.
\end{proof}
\section{Exemplos resolvidos}
Pondo $u=x^2$, temos $du=2x\,dx$:
\[
\int 2x\cos(x^2)\,dx=\int\cos(u)\,du=\sin(x^2)+C.
\]
A derivada da resposta é $2x\cos(x^2)$.
Para $\int xe^x\,dx$, escolha $u=x$ e $dv=e^x\,dx$. Então $du=dx$, $v=e^x$ e
\[
\int xe^x\,dx=xe^x-\int e^x\,dx=(x-1)e^x+C.
\]
Conferência: $\frac{d}{dx}[(x-1)e^x]=xe^x$. Finalmente,
$\int_0^1 xe^x\,dx=[(x-1)e^x]_0^1=1$.
\QuadroEAvaliacao{À esquerda: cadeia/substituição; à direita: produto/partes; na base: verificação por derivadas.}{Por que $\sin(x^2)$ não é primitiva de $\cos(x^2)$? Porque sua derivada é $2x\cos(x^2)$.}
\Auditoria{Não omitir $+C$ em integrais indefinidas; não perder fatores da regra da cadeia; não anunciar aplicações sem resolver o caso definido; distinguir integral indefinida de um número obtido numa integral definida.}
\section{Bibliografia}
J. Stewart, \emph{Cálculo}; E. L. Lima, \emph{Curso de Análise}.

\end{document}
